\documentclass[conference]{IEEEtran}

% Fonts
\usepackage{times}

% Math
\usepackage{amssymb}
\usepackage[cmex10]{amsmath}
\usepackage{amsfonts}
\usepackage{bm}

% Figures and floats (short two-column papers)
\usepackage[hang,center,tight]{subfigure}
\usepackage{float}
\usepackage{graphicx}
\setlength{\textfloatsep}{6pt plus 1pt minus 2pt}
\setlength{\dbltextfloatsep}{6pt plus 1pt minus 2pt}
\setlength{\floatsep}{5pt plus 1pt minus 2pt}
\setlength{\intextsep}{5pt plus 1pt minus 2pt}
\setlength{\abovecaptionskip}{3pt}
\setlength{\belowcaptionskip}{0pt}

% Citations and line breaking
\usepackage[hyphens,obeyspaces]{url}
\usepackage{cite}
\usepackage[none]{hyphenat}
\tolerance=9999
\emergencystretch=10pt
\hyphenpenalty=10000
\exhyphenpenalty=100

% Tables
\usepackage{booktabs}
\usepackage{threeparttable}
\usepackage{multirow}

% Units and symbols (ieee-latex-layout/reference.md §1)
\usepackage{textcomp}
\usepackage{textgreek}

\usepackage[english]{babel}

% Optional: algorithms (uncomment if needed)
% \usepackage{algorithm}
% \usepackage{algorithmicx}
% \usepackage{algpseudocode}
% \renewcommand{\algorithmicrequire}{\textbf{Input:}}
% \renewcommand{\algorithmicensure}{\textbf{Output:}}

\begin{document}
\bstctlcite{IEEEexample:BSTcontrol}
\graphicspath{{./figures/}}

%----------------------------------------------------------------------
% Title, abstract, keywords
%----------------------------------------------------------------------

\title{Paper Title}

\author{
Author~One$^{1}$,
Author~Two$^{1}$, and
Author~Three\textsuperscript{1,$\ddagger$} \\
$^{1}$Affiliation, City, Country \\
$\ddagger$Email: author@institution.edu
\thanks{Acknowledgment text (funding, corresponding author).}
}

\IEEEoverridecommandlockouts
\maketitle

\begin{abstract}
State the problem, method, and strongest quantitative evidence under the adopted
setup. Use the same numbers as the results section.
\end{abstract}

\begin{IEEEkeywords}
Keyword one, keyword two, keyword three, keyword four, keyword five.
\end{IEEEkeywords}

%----------------------------------------------------------------------
\section{Introduction}
\label{sec:introduction}
%----------------------------------------------------------------------

\IEEEPARstart{F}{irst} sentence: application motivation.
Gap, prior limitations, proposed approach, and numbered contributions.

%----------------------------------------------------------------------
\section{Methodology}
\label{sec:methodology}
%----------------------------------------------------------------------

Objective methodology wording: what is adopted, held fixed, measured, compared.

%----------------------------------------------------------------------
\section{Results}
\label{sec:results}
%----------------------------------------------------------------------

% Example single-column figure
% \begin{figure}[!t]
%     \centering
%     \includegraphics[width=\columnwidth]{figures/fig_example.pdf}
%     \caption{Concise takeaway and scope.}
%     \label{fig:example}
% \end{figure}

% Example full-width figure
% \begin{figure*}[!t]
%     \centering
%     \includegraphics[width=0.76\textwidth]{figures/fig_wide.pdf}
%     \caption{Concise takeaway and scope.}
%     \label{fig:wide}
% \end{figure*}

%----------------------------------------------------------------------
\section{Conclusion}
\label{sec:conclusion}
%----------------------------------------------------------------------

Summarize what was proposed, demonstrated, enabled, and what scope remains bounded.

%----------------------------------------------------------------------
% References (compact font for short papers)
%----------------------------------------------------------------------
\begingroup
\makeatletter
\let\oldthebibliography\thebibliography
\let\endoldthebibliography\endthebibliography
\renewenvironment{thebibliography}[1]{%
    \oldthebibliography{#1}%
    \fontsize{7}{7.15}\selectfont
    \setlength{\itemsep}{0pt}%
    \setlength{\parskip}{0pt}%
}{\endoldthebibliography}
\makeatother
\bibliographystyle{IEEEtran}
\bibliography{./IEEEabrv,./reference}
\endgroup

\end{document}
